<preamble>

% Log to stderr as well as the log file.
\tracingonline=1

% Don't stop when messages are printed to the terminal
\nonstopmode

% Output up to 1 million nodes.
\showboxbreadth=1000000
% Output up to 100 nested boxes.
\showboxdepth=100

% Make the page really wide and then skip the first pass of the line-breaking algorithm.
% This forces the hyphenation to kick in.
\pretolerance=-1
\hsize=16383pt

% Don't output anything to the DVI file.
\output={
    % This is needed to avoid warnings about some items not being output.
    \setbox254=\box255
    \setbox255=\hbox{} \shipout\box255
}

% Macro that prints the contents on its own line in the terminal.
\def\fullLineMessage#1{%
    {%
        \newlinechar=`@%
        \message{@#1@}%
    }%
}

\def\buildAndPrintBoxes#1{

% Flush the output of the main vertical list so that any content already there
% won't be seen in this macro expansion.
\eject

% (1) First we print a regular hbox containing the text.
% The only catch here is that the text is set in internal horizontal mode,
% which is a little different than regular horiztonal mode. Specfically, some
% manual discretionaries are not created.
\fullLineMessage{internal horizontal box start}
\setbox253=\hbox{#1}
\showbox253
\fullLineMessage{internal horiztonal box end}

% (2) Next print the text as it is put on the main horizontal list.
\noindent
#1%
\fullLineMessage{unhyphenated hlist start}\showlists\fullLineMessage{unhyphenated hlist end}

% (3) After the last paragraph ended, TeX ran the line breaking algoirthm, performed
% hyphenation (becaues \pretolerance=-1) and put the resulting hbox on the vertical list
% (a single hbox because the page is so wide). So printing the current lists will
% show this hbox.
\fullLineMessage{hyphenated hlist start}\showlists\fullLineMessage{hyphenated hlist end}

}

<print_calls>

\end
